\beginchapter Chapter 19. Displayed Equations

By now you know how to type mathematical formulas so that \TeX\ will handle
them with supreme elegance; your knowledge of math typing is nearly complete.
But there is one more part to the story, and the purpose of this chapter
is to present the happy ending. We have discussed how to deal with individual
formulas; but ^{displays} often involve a whole bunch of different formulas,
or different pieces of a huge formula, and it's a bit of a problem to
lay them out so that they line up properly with each other. Fortunately,
large displays generally fall into a few simple patterns.

\subsection One-line displays. Before plunging into the general question
of display layout, let's recapitulate what we have already covered.  If
you type `|$$|\<formula>|$$|', \TeX\ will display the formula
^^{dollardollar} in flamboyant display style, centering it on a line by
itself. We have also noted in Chapter~18 that it's possible to display two
short formulas at once, by typing
`|$$|\<formula$_1$>^|\qquad|\<formula$_2$>|$$|'; this reduces the
two-formula problem to a one-formula problem. You get the two formulas
separated by two quads of space, the whole being centered on a line.

Displayed equations often involve ordinary text. Chapter~18 explains
how to get roman type into formulas without leaving math mode, but the
best way to get text into a display is to put it into an ^|\hbox|.
There needn't even be any math at all; to typeset
$$\hbox{Displayed Text}$$
you can simply say `|$$\hbox{Displayed Text}$$|'. But here's a more interesting
example:
$$X_n=X_k \qquad\hbox{if and only if}\qquad
    Y_n=Y_k \quad\hbox{and}\quad Z_n=Z_k.$$
Formulas and text were combined in this case by typing
\begintt
$$X_n=X_k \qquad\hbox{if and only if}\qquad
    Y_n=Y_k \quad\hbox{and}\quad Z_n=Z_k.$$
\endtt
Notice that |\qquad| appears around `if and only if', but a single ^|\quad|
surrounds `and'; this helps to indicate that the $Y$ and~$Z$ parts of the
display are related more closely to each other than to the $X$~part.

Consider now the display
$$Y_n=X_n\bmod p \quad\hbox{and}\quad Z_n=X_n\bmod q
    \qquad\hbox{for all }n\ge0.$$
Can you figure out how to type this? One solution is
\begintt
$$Y_n=X_n\bmod p \quad\hbox{and}\quad Z_n=X_n\bmod q
    \qquad\hbox{for all }n\ge0.$$
\endtt
Notice that a space has been left after `|all|' in the hbox here, since spaces
disappear when they are out in formula-land. But there's a simpler and more
logical way to proceed, once you get used to \TeX's idea of modes: You can type
\begintt
... \qquad\hbox{for all $n\ge0$.}$$
\endtt
Wow---that's math mode inside of horizontal mode inside of display
math mode. But in this way your manuscript mirrors what you are trying to
accomplish, while the previous solution (with the space after `|all|')
looks somewhat forced.

\exercise Typeset the following four displays (one at a time):
$$\openup1\jot\displaylines{
\sum_{n=0}^\infty a_nz^n\qquad\hbox{converges if}\qquad
  \vert z\vert<\Bigl(\limsup_{n\to\infty}
    \root n\!\of{\vert a_n\vert}\,\Bigr)^{-1}.\cr
{f(x+\Delta x)-f(x)\over\Delta x}\to f'(x)\qquad\hbox{as $\Delta x\to0$.}\cr
\noalign{\vskip2pt}
\Vert u_i\Vert=1,\qquad u_i\cdot u_j=0\quad\hbox{if $i\ne j$.}\cr
\it\hbox{The confluent image of}\quad\left\{
  \matrix{\hbox{an arc}\hfill\cr\hbox{a circle}\hfill\cr
    \hbox{a fan}\hfill\cr}
  \right\}\quad\hbox{is}\quad\left\{
  \matrix{\hbox{an arc}\hfill\cr\hbox{an arc or a circle}\hfill\cr
    \hbox{a fan or an arc}\hfill\cr}\right\}.\cr
\noalign{\vskip-8pt}}$$
^^|\Delta| ^^|\Vert|
% the last example comes from Proc AMS 55 (1976), 410, with typos corrected
\answer |$$\sum_{n=0}^\infty a_nz^n\qquad\hbox{converges if}\qquad|\parbreak
|  |\||z|\||<\Bigl(\limsup_{n\to\infty}\root n\!\of{|\||a_n|\|^^|\root|
  |}\,\Bigr)^{-1}.$$|\kern-.33pt\par
\smallskip
|$${f(x+\Delta x)-f(x)\over\Delta x}\to f'(x)|\parbreak
|    \qquad\hbox{as $\Delta x\to0$.}$$|\par
\smallskip
|$$\|\||u_i\|\||=1,\qquad u_i\cdot u_j=0\quad\hbox{if $i\ne j$.}$$|\par
\smallskip
|$$\it\hbox{The confluent image of}\quad\left\{|\parbreak
|    \matrix{\hbox{an arc}\hfill\cr\hbox{a circle}\hfill\cr|\parbreak
|      \hbox{a fan}\hfill\cr}|\parbreak
|    \right\}\quad\hbox{is}\quad\left\{|\parbreak
|    \matrix{\hbox{an arc}\hfill\cr|\parbreak
|      \hbox{an arc or a circle}\hfill\cr|\parbreak
|      \hbox{a fan or an arc}\hfill\cr}\right\}.$$|\par
\smallskip\noindent
The first example includes |\!| and |\,| to give slightly refined spacing;
but the point of the problem was to illustrate the hbox, not to fuss over
such extra details.
The last example can be done much more simply using the ideas of
Chapter~22, if you don't mind descending to the level of \TeX\ primitives;
for example, the first matrix could be replaced by ^^|\halign|
\begintt
\,\vcenter{\halign{#\hfil\cr an arc\cr a circle\cr a fan\cr}}\,
\endtt
and the second is similar.

\dangerexercise Sometimes display style is too grandiose, when the formula
being displayed is
$$y={1\over2}x$$
or something equally simple. One day B. L. ^{User} tried to remedy this by
typing it as `|$$y={\scriptstyle1\over\scriptstyle2}x$$|', but the
resulting formula
$$y={\scriptstyle1\over\scriptstyle2}x$$
wasn't at all what he had in mind. What's the right way to get ^^{one half}
^^{1/2--unslashed form} simply `$y={1\over2}x$' when you don't want big
^{fractions in displays}?
\answer |$$\textstyle y={1\over2}x$$|. \ (Switching to text style is
especially common in multiline formulas. For example, you will probably
find occasions to use ^|\textstyle| on both sides of the |&|'s within
an ^|\eqalign|.)

\dangerexercise What difference, if any, is there between the result of
typing `|$$|\<formula>|$$|' and the result of typing
`|$$\hbox{$|\<formula>|$}$$|'\thinspace?
\answer The latter formula will be in text style, not display style.
And even if you do type `|$$\hbox{$\displaystyle{|\<formula>|}$}$$|', the
results are not quite the same, as we will see later: \TeX\ will compress
the glue in `|$$|\<formula>|$$|' if the formula is too wide to fit on
a line at its natural width, but the glue inside |\hbox{...}| is frozen
at its natural width. The |\hbox| version also invokes
|\everyhbox| and |\everymath|.

\dangerexercise You may have noticed that most of the displays in this
manual are not centered; displayed material is usually aligned at the
left with the paragraph indentation, as part of the book design, because
this is an unusual book. Explain how you could typeset a formula like
$$\leftline{\indent$\displaystyle
    1-{1\over2}+{1\over3}-{1\over4}+\cdots=\ln2$}$$
^^{displays, non-centered}
that is off-center in this way.
\answer One solution is to put the formula in an hbox that occupies a full line:
\begintt
$$\leftline{\indent$\displaystyle
    1-{1\over2}+{1\over3}-{1\over4}+\cdots=\ln2$}$$
\endtt
But this takes a bit of typing. If you make the definitions
\begintt
\def\leftdisplay#1$${\leftline{\indent$\displaystyle{#1}$}$$}
\everydisplay{\leftdisplay}
\endtt
you can type `|$$|\<formula>|$$|' as usual, and the formatting will be
inserted automatically. \ (This doesn't work with equation numbers;
Appendix~D illustrates how to handle them as well.)

If you've had previous experience typing mathematical papers, you probably
have been thinking, ``What about ^{equation numbers}? When is this book
going to talk about them?'' Ah yes, now is the time to discuss those sneaky
little labels that appear off to the side of displays. If you type
\begindisplay
|$$|\<formula>|\eqno|\<formula>|$$|
\enddisplay
\TeX\ will display the first formula and it will also put an equation number
(the second formula) at the right-hand margin. For example,
\begintt
$$x^2-y^2 = (x+y)(x-y).\eqno(15)$$
\endtt
^^|\eqno|
will produce this:
$$x^2-y^2 = (x+y)(x-y).\eqno(15)$$
You can also get equation numbers at the left-hand margin, with ^|\leqno|.
For example,
\begintt
$$x^2-y^2 = (x+y)(x-y).\leqno(16)$$
\endtt
will produce this:
$$x^2-y^2 = (x+y)(x-y).\leqno(16)$$
Notice that you always give the equation number second, even when it is going
to appear at the left. Everything from the |\eqno| or |\leqno| command to
the |$$| that ends the display is the equation number. Thus, you're not
allowed to have two equation numbers in the same display; but there's a
way to get around that restriction, as we'll see later.

\danger Nowadays people are using right-hand equation numbers more and more,
because a display most often comes at the end of a sentence or clause, and
the right-hand convention keeps the number from intruding into the clause.
Furthermore, it's often possible to save space when a displayed equation
follows a short text line, since less space is needed above the display;
such savings are not possible with |\leqno|, because there's no room for
overlap. For example, there is less space above display~(15) than there is
above~(16) in our illustrations of\/ |\eqno| and |\leqno|, although the
formulas and text are otherwise identical.

\danger If you look closely at (15) and (16) above, you can see that the
displayed formulas have been centered without regard to the presence of
the equation numbers. But when a formula is large, \TeX\ makes sure that
it does not interfere with its number; the equation number may even be
placed on a line by itself.

\exercise How would you produce the following display?
$$\prod_{k\ge0}{1\over(1-q^kz)}=
    \sum_{n\ge0}z^n\bigg/\!\!\prod_{1\le k\le n}(1-q^k).\eqno(16')$$
\answer |$$\prod_{k\ge0}{1\over(1-q^kz)}=|\parbreak
        |   \sum_{n\ge0}z^n\bigg/\!\!\prod_{1\le k\le n}(1-q^k).\eqno(16')$$|

\dangerexercise Equation numbers are math formulas, typeset in text style.
So how can you get an equation number like `\hbox{(3--1)}'
(with an ^{en-dash})?
\answer |\eqno\hbox{(3--1)}|.

\ddangerexercise B. L. ^{User} tried typing `|\eqno(*)|' and `|\eqno(**)|',
and he was pleased to discover that this produced the equation numbers
`$(*)$' and `$(**)$'. \ [He had been a bit worried that they would come out
`(*)' and `(**)' instead.] \ But then a few months later he tried
`|\eqno(***)|' and got a surprise. What was it?
\answer When you type an ^{asterisk} in math mode, plain \TeX\ considers
|*| to be a binary operation. In the cases `|(*)|' and `|(**)|', the
binary operations are converted to type~Ord, because they don't appear in
a binary context; but the middle asterisk in `|(***)|' remains of type~Bin.
So the result was `$(***)$'. To avoid the extra medium spaces, you can
type `|\eqno(*{*}*)|'; or you can change ^|\mathcode||`*|, if you never use
|*| as a binary operation.

\ddanger Somewhere in this manual there ought to be a description of exactly
how \TeX\ displays formulas; i.e., how it centers them, how it places the
equation numbers, how it inserts extra space above and below, and so on.
Well, now is the time for those rules to be stated. They are somewhat
complex, because they interact with things like |\parshape|, and because
they involve several parameters that haven't been discussed yet. The purpose
of the rules is to explain exactly what sorts of boxes, glue, and
penalties are placed onto the current ^{vertical list} when a display occurs.

\ddanger If a display occurs after, say, four lines of a paragraph, \TeX's
internal register called ^|\prevgraf| will be equal to~4 when the display
starts.  The display will be assumed to take three lines, so |\prevgraf|
will become~7 when the paragraph is resumed at the end of the display
(unless you have changed |\prevgraf| in the meantime). \TeX\ assigns
special values to three \<dimen> parameters immediately after the opening
|$$| is sensed:  ^|\displaywidth| and ^|\displayindent| are set to the
line width~$z$ and the shift amount~$s$ for line number |\prevgraf|${}+2$,
based on the current paragraph shape or hanging indentation. \ (Usually
|\displaywidth| is the same as ^|\hsize|, and |\displayindent| is zero,
but the paragraph shape can vary as described in Chapter~14.) \ Furthermore,
^|\predisplaysize| is set to the effective width~$p$ of the line preceding
the display, as follows: If there was no previous line (e.g., if the |$$|
was preceded by ^|\noindent| or by the closing |$$| of another display),
$p$~is set to $-16383.99999\pt$ (i.e., to the smallest legal dimension,
$-$^|\maxdimen|).  Otherwise \TeX\ looks inside the hbox that was formed
by the previous line, and sets $p$ to the position of the right edge of
the rightmost box inside that hbox, plus the indentation by which the
enclosing hbox has been moved right, plus two ems in the current font.
However, if this value of~$p$ depends on the fact that glue in that hbox
was stretching or shrinking---for example, if the ^|\parfillskip| glue is
finite, so that the material preceding it has not been set at its natural
width---then $p$~is set to |\maxdimen|. \ (This doesn't happen
often, but it keeps \TeX\ machine independent, since $p$~never depends on
quantities that may be rounded differently on different computers.)  \
Notice that |\displaywidth| and |\displayindent| are not
affected by |\leftskip| and |\rightskip|, but
|\predisplaysize| is. The values of\/ |\displaywidth|,
|\displayindent|, and |\predisplaysize|
will be used by \TeX\ after the displayed formula has been
read, as explained below; your program can examine them and/or change
them, if you want the typesetting to be done differently.

\ddanger After a display has been read, \TeX\ converts it from a math list
to a horizontal list~$h$ in display style, as explained in
Appendix~G\null.  An equation number, if present, is processed in text
style and put into an hbox~$a$ with its natural width. Now the fussy
processing begins: Let $z$, $s$, and~$p$ be the current values of\/\
|\displaywidth|, |\displayindent|, and
|\predisplaysize|. Let $q$ and~$e$ be zero if there is no equation
number; otherwise let~$e$ be the width of the equation number, and let~$q$
be equal to $e$~plus one quad in the symbols font (i.e., in
^|\textfont||2|). Let $w_0$ be the natural width of the displayed
formula~$h$. If $w_0+q\le z$, list~$h$~is packaged in an hbox~$b$ having
its natural width~$w_0$.  But if $w_0+q>z$ (i.e., if the display is too
wide to fit at its natural width), \TeX\ performs the following ``^{squeeze
routine}'':  If $e\ne0$ and if there is enough shrinkability in the
displayed formula~$h$ to reduce its width to $z-q$, then list $h$ is
packaged in an hbox~$b$ of width~$z-q$. Otherwise $e$ is set to zero, and
list~$h$ is packaged in a (possibly overfull) hbox~$b$ of width $\min(w_0,z)$.

\ddanger (Continuation.) \ \TeX\ tries now to center the display without
regard to the equation number. But if such centering would make it too close
to that number (where ``too close'' means that the space between them is less
than the width~$e$), the equation is either centered in the remaining space
or placed as far from the equation number as possible. The latter alternative
is chosen only if the first item on list~$h$ is glue, since \TeX\ assumes that
such glue was placed there in order to control the spacing precisely.
But let's state the rules more formally: Let~$w$ be the width of box~$b$.
\TeX\ computes a displacement~$d$, to be used later when positioning box~$b$,
by first setting $d={1\over2}(z-w)$. If $e>0$ and if $d<2e$, then $d$~is
reset to ${1\over2}(z-w-e)$ or to zero, where zero is chosen if list~$h$
begins with a glue item.

\ddanger (Continuation.) \ \TeX\ is now ready to put things onto the current
vertical list, just after the material previously constructed for the
paragraph-so-far. First comes a ^{penalty} item, whose cost is an integer
parameter called ^|\predisplaypenalty|. Then comes glue. If $d+s\le p$,
or if there was a left equation number (|\leqno|), \TeX\ sets $g_a$ and
$g_b$ to glue items specified by the parameters ^|\abovedisplayskip|
and ^|\belowdisplayskip|, respectively; otherwise $g_a$ and $g_b$ become
glue items corresponding to ^|\abovedisplayshortskip| and
^|\belowdisplayshortskip|. \ [Translation: If the predisplaysize is short
enough so that it doesn't overlap the displayed formula, the glue above and
below the display will be ``short'' by comparison with the glue that is
used when there is an overlap.] \ If $e=0$ and if there is an |\leqno|,
the equation number is appended as an hbox by itself, shifted right~$s$ and
preceded by interline glue as usual; an infinite penalty is also appended,
to prevent a page break between this number and the display. Otherwise
a glue item~$g_a$ is placed on the vertical list.

\ddanger (Continuation.) \ Now comes the displayed equation itself. If
$e\ne0$, the equation number box~$a$ is combined with the formula box~$b$ as
follows: Let~$k$ be a kern of width $z-w-e-d$. In the |\eqno| case, box~$b$
is replaced by an hbox containing $(b,k,a)$; in the |\leqno| case, box~$b$
is replaced by an hbox containing $(a,k,b)$, and $d$~is set to zero. In all
cases, box~$b$ is then appended to the vertical list, shifted right by~$s+d$.

\ddanger (Continuation.) \ The final task is to append the glue or the
equation number that follows the display. If there was an |\eqno| and if
$e=0$, an infinite penalty is placed on the vertical list, followed by the
equation number box~$a$ shifted right by $s+z$ minus its width, followed
by a penalty item whose cost is the value of\/ ^|\postdisplaypenalty|.
Otherwise a penalty item for the |\postdisplaypenalty| is appended
first, followed by a glue item for~$g_b$ as specified above. \TeX\ now
adds~3 to |\prevgraf| and returns to horizontal mode, ready to resume the
paragraph.

\ddanger One consequence of these rules is that you can force an equation
number to appear on a line by itself by making its width zero, i.e.,
by saying either `|\eqno|^|\llap||{$|\<formula>|$}|' or
`|\leqno|^|\rlap||{$|\<formula>|$}|'. This makes $e=0$, and
the condition $e=0$ controls \TeX's positioning logic, as explained
in the rules just given.

\ddanger Plain \TeX\ sets |\predisplaypenalty=10000|, because fine
printers traditionally shun displayed formulas at the very top of a page.
You can change |\predisplaypenalty| and
|\postdisplaypenalty| if you want to encourage or discourage
page breaks just before or just after a display. For example,
`\hbox{|$$\postdisplaypenalty=|}\allowbreak\hbox{|-10000|\<formula>|$$|}'
will force a page break, putting the formula at the bottom line. It is better
to force a ^{page break} this way than to say ^|\eject| right after |$$...$$|;
such an eject (which follows the |\belowdisplayskip| glue below the
display) causes the page to be short, because it leaves unwanted glue at
the bottom.

\ddangerexercise Read the rules carefully and deduce the final position
of `$x=y$' in the formula
\begintt
$$\quad x=y \hskip10000pt minus 1fil \eqno(5)$$
\endtt
assuming that there is no hanging indentation. Also consider |\leqno|
instead of\/ |\eqno|.
\answer Assuming that |\hsize| is less than $10000\pt$, the natural width of
this equation will be too large to fit on a line; also, |\quad| specifies
glue at the left. Therefore `$x=y$' will appear exactly $1\em$ from
the left, and `(5)' will appear flush right. \ (The widths will satisfy
^^{displays, non-centered} $w=z-q$, $d=0$, $k=q-e=18\rm\,mu$.) \
In the case of\/ |\leqno|, `(5)' will appear flush left, followed by
one quad of space in |\textfont2|, followed by one quad of space in the
current text font, followed by `$x=y$'.

\ddanger \TeX\ also allows ``^{alignment displays},'' which are not
processed in math mode because they contain no formulas at the outer
level. An alignment display is created by commands of the general form
\begindisplay
|$$|\<assignments>|\halign{|\<alignment>|}|\<assignments>|$$|
\enddisplay
where the \<assignments> are optional things like parameter changes that do not
produce any math lists. In such displays, the |\halign| is processed exactly
as if it had appeared in vertical mode, and it will construct a vertical
list~$v$ as usual, except that each row of the alignment will be shifted
right by the ^|\displayindent|. After the alignment and the closing
assignments have been processed, \TeX\ will put a ^|\predisplaypenalty|
item and some ^|\abovedisplayskip| glue on the main vertical list, followed
by~$v$, followed by a ^|\postdisplaypenalty| item and ^|\belowdisplayskip|
glue. Thus, alignment displays are essentially like ordinary alignments,
except that they can interrupt paragraphs; furthermore, they are embedded in
glue and penalties just like other displays. The ^|\displaywidth| and
^|\predisplaysize| do not affect the result, although you could use
those parameters in your ^|\halign|. An entire alignment display is considered
to be only three lines long, as far as ^|\prevgraf| is concerned.

\subsection Multi-line displays. OK, the use of displayed formulas is
very nice. But when you try typing a lot of manuscripts you will run into
some displays that don't fit the simple pattern of a one-line formula with
or without an equation number. Plain \TeX\ provides special control
sequences that will cover most of the remaining cases.

Multi-line displays usually consist of several equations that should be
lined up by their `$=$'~signs, as in
$$\eqalign{X_1+\cdots+X_p&=m,\cr
           Y_1+\cdots+Y_q&=n.\cr}$$
The recommended procedure for such a display is to use ^|\eqalign|,
which works with special markers |&| ^^{ampersand} and ^|\cr| that we
have already encountered in connection with |\cases| and |\matrix|
in Chapter~18. Here's how to type this particular one:
\begintt
$$\eqalign{X_1+\cdots+X_p&=m,\cr
           Y_1+\cdots+Y_q&=n.\cr}$$
\endtt
There can be any number of equations in an |\eqalign|; the general pattern is
$$\halign{\indent#&#\hfil\cr
|\eqalign{|&\<left-hand side$_1$>|&|\<right-hand side$_1$>|\cr|\cr
           &\<left-hand side$_2$>|&|\<right-hand side$_2$>|\cr|\cr
\noalign{\vskip-2pt}
           &\qquad\vdots\cr
           &\<left-hand side$_n$>|&|\<right-hand side$_n$>|\cr}|\cr}$$
where each \<right-hand side> starts with the symbol on which you want
alignment to occur. For example, every right-hand side often begins
with an $=$~sign. The equations will be typeset in display style.

\exercise In practice, the left-hand sides of aligned formulas are often
blank, and the alignment is often done with respect to other symbols
as well as~$=$. For example, the following display is typical; see if you
can guess how the author typed it:
$$\eqalign{T(n)\le T(2^{\lceil\lg n\rceil})
    &\le c(3^{\lceil\lg n\rceil}-2^{\lceil\lg n\rceil})\cr
    &<3c\cdot3^{\lg n}\cr
    &=3c\,n^{\lg3}.\cr}$$ % from v2 p279
\answer (Note in particular that the final `|.|'\ comes {\sl before\/} the
final `|\cr|'.)
\begintt
$$\eqalign{T(n)\le T(2^{\lceil\lg n\rceil})
    &\le c(3^{\lceil\lg n\rceil}-2^{\lceil\lg n\rceil})\cr
    &<3c\cdot3^{\lg n}\cr
    &=3c\,n^{\lg3}.\cr}$$
\endtt

The result of\/ |\eqalign| is a vertically centered box. This makes it easy to
get a formula like
$$\left\{
\eqalign{\alpha&=f(z)\cr \beta&=f(z^2)\cr \gamma&=f(z^3)\cr}
\right\}\qquad\left\{
\eqalign{x&=\alpha^2-\beta\cr y&=2\gamma\cr}\right\}.$$ % meaningless
You simply use |\eqalign| twice in the same line:
\begintt
$$\left\{
\eqalign{\alpha&=f(z)\cr \beta&=f(z^2)\cr \gamma&=f(z^3)\cr}
\right\}\qquad\left\{
\eqalign{x&=\alpha^2-\beta\cr y&=2\gamma\cr}\right\}.$$
\endtt

\exercise Try your hand at the numbered two-line display % Polya/Szego V.29
$$\eqalign{P(x)&=a_0+a_1x+a_2x^2+\cdots+a_nx^n,\cr
   P(-x)&=a_0-a_1x+a_2x^2-\cdots+(-1)^na_nx^n.\cr}\eqno(30)$$
[{\sl Hint:\/} Use the fact that |\eqalign| produces a vertically centered
box; the equation number `(30)' is supposed to appear halfway between
the two lines.]
\answer |$$\eqalign{P(x)&=a_0+a_1x+a_2x^2+\cdots+a_nx^n,\cr|\parbreak
        |   P(-x)&=a_0-a_1x+a_2x^2-\cdots+(-1)^na_nx^n.\cr}\eqno(30)$$|\par

\exercise What happens if you forget the |&| in one equation of an |\eqalign|?
\answer Both sides of that equation are considered to be on the left, so
you get results that look like this:
$$\openup-\jot
\left\{\eqalign{\alpha&=f(z)\cr \beta&=f(z^2)\cr \gamma=f(z^3)\cr}
  \right\}.$$

\danger Multi-line formulas sometimes fit together in odd ways, and you'll
find that every once in a~while you will want to move certain lines farther
apart or closer together. If you type `^|\noalign||{|^|\vskip|\<glue>|}|'
after any |\cr|, \TeX\ will insert the given amount of extra glue just
after that particular line. For example,
\begintt
\noalign{\vskip3pt}
\endtt
will put $3\pt$ of additional space between lines. You can also change the
amount of space before the first line, in the same way.

The next level of complexity occurs when you have several aligned
equations with several equation numbers. Or perhaps some of the
lines are numbered and others are not:
$$\eqalignno{(x+y)(x-y)&=x^2-xy+yx-y^2\cr
    &=x^2-y^2;&(4)\cr
  (x+y)^2&=x^2+2xy+y^2.&(5)\cr}$$
For this situation plain \TeX\ provides ^|\eqalignno|; you use it like
|\eqalign|, but on each line that you want an equation number you add
`|&|\<equation number>' just before the |\cr|. The example above was
generated by
\begintt
$$\eqalignno{(x+y)(x-y)&=x^2-xy+yx-y^2\cr
    &=x^2-y^2;&(4)\cr
  (x+y)^2&=x^2+2xy+y^2.&(5)\cr}$$
\endtt
Notice that the second |&| is omitted unless there's an equation number.

And there's also ^|\leqalignno|, which puts equation numbers at the left.
In this case it is appropriate to move the `(4)' to the beginning
of its equation:
$$\leqalignno{(x+y)(x-y)&=x^2-xy+yx-y^2&(4)\cr
    &=x^2-y^2;\cr
  (x+y)^2&=x^2+2xy+y^2.&(5)\cr}$$
Although the equation numbers appear at the left, you are still supposed to
input them at the right, just as you do with |\leqno|; in other words,
you should type
`|$$\leqalignno{(x+y)(x-y)&...&(4)\cr...}$$|' to get the previous display.

Caution: |\eqalignno| and |\leqalignno| both center the set of equations
without regard to the widths of the equation numbers. If the equations or
their numbers get too wide, they might overlap, yet no error message will
be given.

\exercise Typeset the following display: ^^|\gcd|
$$\leqalignno{\gcd(u,v)&=\gcd(v,u);&(9)\cr
    \gcd(u,v)&=\gcd(-u,v).&(10)\cr}$$ % v2 p316
\answer |$$\leqalignno{\gcd(u,v)&=\gcd(v,u);&(9)\cr|\parbreak
        |    \gcd(u,v)&=\gcd(-u,v).&(10)\cr}$$|

\exercise And here's another one to try, just to keep in practice: ^^|\int|
$$\vbox{
\eqalignno{\biggl(\int_{-\infty}^\infty e^{-x^2}\,dx\biggr)^2
  &=\int_{-\infty}^\infty\int_{-\infty}^\infty
    e^{-(x^2+y^2)}\,dx\,dy\cr
  &=\int_0^{2\pi}\int_0^\infty e^{-r^2}r\,dr\,d\theta\cr
  &=\int_0^{2\pi}\biggl(-{e^{-r^2}\over2}
    \bigg\vert_{r=0}^{r=\infty}\,\biggr)\,d\theta\cr
  &=\pi.&(11)\cr}
}$$ % cf Joy of TeX
\answer %
|$$\eqalignno{\biggl(\int_{-\infty}^\infty e^{-x^2}\,dx\biggr)^2|\parbreak
        |  &=\int_{-\infty}^\infty\int_{-\infty}^\infty|\parbreak
        |    e^{-(x^2+y^2)}\,dx\,dy\cr|\parbreak
        |  &=\int_0^{2\pi}\int_0^\infty e^{-r^2}r\,dr\,d\theta\cr|\parbreak
        |  &=\int_0^{2\pi}\biggl(-{e^{-r^2}\over2}|\parbreak
        |    \bigg|\||_{r=0}^{r=\infty}\,\biggr)\,d\theta\cr|\parbreak
        |  &=\pi.&(11)\cr}$$| ^^|\bigg|

\danger Although |\eqalign| and |\eqalignno| look nearly the same, there's
really a fundamental distinction between them: |\eqalign| makes a single,
vertically centered box, which is no wider than it needs to be; but
|\eqalignno| generates a set of lines that have the full display width
(reaching all the way to both margins). Thus, for example, you can use
|\eqalign| several times in a display, but |\eqalignno| can appear only
once. If you try to use ^|\eqno| in conjunction with |\eqalign|,
you get a decent result, but if you try to use |\eqno| in connection
with |\eqalignno| you'll get some sort of weird error message(s).

\ddanger The definitions in Appendix~B reveal why |\eqalign| and |\eqalignno|
behave differently: |\eqalign| is an
abbreviation for ^|\vcenter||{|^|\halign||{...}}|, while
|\eqalignno| is an abbreviation for |\halign to\displaywidth{...}|;
thus the |\eqalignno| macro generates an ``^{alignment display}.''

\ddanger This difference between |\eqalign| and |\eqalignno| has two
interesting consequences. \ (1)~It's impossible to break an |\eqalign|
between pages, but an |\eqalignno| can be broken. In fact, you can
{\sl force\/} a ^{page break} after a particular line if you insert
`^|\noalign||{|^|\break||}|'
after the |\cr| for that line. You can prohibit {\sl all\/} breaks
in an |\eqalignno| if you set ^|\interdisplaylinepenalty||=10000|; or you
can enclose the whole works in a ^|\vbox|:
\begintt
$$\vbox{\eqalignno{...}}$$
\endtt
(2) You can also insert a line of text between two equations, without
losing the alignment. For example, consider the two displays
$$\eqalignno{x&=y+z\cr
  \noalign{\hbox{and}}
  x^2&=y^2+z^2.\cr}$$
These were actually generated as a single display by typing
\begintt
$$\eqalignno{x&=y+z\cr
  \noalign{\hbox{and}}
  x^2&=y^2+z^2.\cr}$$
\endtt
Therefore the fact that their $=$ signs line up is not just a lucky
coincidence. Sometimes you will want to adjust the spacing above or below
such a line of inserted text, by putting a |\vskip| or two inside of the
|\noalign{...}|.  Incidentally, this example also shows that it is
possible to use |\eqalignno| without giving any equation numbers.

\ddangerexercise What happens if\/ |\eqalign| is substituted for
|\eqalignno| in this last example?
\answer You get the displayed box
$$\eqalign{x&=y+z\cr
  \noalign{\hbox{and}}
  x^2&=y^2+z^2.\cr}$$
Reason: The `and' occurs at the left of the |\eqalign| box, not at the
left of the whole display, and the |\eqalign| box is centered as usual.

\ddangerexercise Our friend Ben ^{User} got into trouble again when he tried to
move an equation number up higher than its usual position, by typing this:
^^|\raise|
\begintt
$$\eqalignno{...&\raise6pt\hbox{(5)}\cr}$$
\endtt
What was his oversight, and what could he have done instead?
\answer By raising the equation number, he increased the line height,
so \TeX\ put extra space between that line and the previous line
when it calculated the interline glue. If he had said
`^|\smash||{\raise...}|', he wouldn't have had that problem.

\danger For other types of displays, plain \TeX\ provides ^|\displaylines|,
which lets you display any number of formulas in any way you want,
without any alignment. The general form is
$$\halign{\indent\hfil#&#\hfil\cr
|$$\displaylines{|&\<displayed formula$_1$>|\cr|\cr
    &\<displayed formula$_2$>|\cr|\cr
\noalign{\vskip-2pt}
    &\qquad\vdots\cr
    &\<displayed formula$_n$>|\cr}$$|\cr}$$
Each formula will be centered, because |\displaylines| puts ^|\hfil| at
the left and the right of each line; you can override this centering to
get things flush left or flush right by inserting ^|\hfill|, which takes
precedence over |\hfil|.

\dangerexercise Use |\displaylines| to typeset the three-line display
$$\displaylines{\hfill x\equiv x;\hfill\llap{(1)}\cr
   \hfill\hbox{if}\quad x\equiv y\quad\hbox{then}\quad
      y\equiv x;\hfill\llap{(2)}\cr
   \hfill\hbox{if}\quad x\equiv y\quad\hbox{and}\quad
      y\equiv z\quad\hbox{then}\quad
      x\equiv z.\hfill\llap{(3)}\cr}$$
\answer |$$\displaylines{\hfill x\equiv x;\hfill\llap{(1)}\cr|\parbreak
        |   \hfill\hbox{if}\quad x\equiv y\quad\hbox{then}\quad|\parbreak
        |      y\equiv x;\hfill\llap{(2)}\cr|\parbreak
        |   \hfill\hbox{if}\quad x\equiv y\quad\hbox{and}\quad|\parbreak
        |      y\equiv z\quad\hbox{then}\quad|\parbreak
        |      x\equiv z.\hfill\llap{(3)}\cr}$$|\par\medskip\noindent
There's also a trickier solution, which begins with
\begintt
$$\displaylines{x\equiv x;\hfil\llap{(1)}\hfilneg\cr
\endtt

\danger If you look closely at the multi-line displays in this chapter,
you'll see that the baselines are farther apart than they are in normal
text; mathematics publishers generally do this in order to make the
displays easier to read. In accordance with this tradition, |\eqalign| and
its relatives automatically increase the ^|\baselineskip|. If~you are
making a multi-line display with \TeX's primitive ^|\halign| command,
instead of using one of the plain \TeX\ macros, you might want to make
this same baseline adjustment, and you can do it easily by saying
`|$$\openup1\jot \halign{...}$$|'. The ^|\openup| macro increases
^|\lineskip| and ^|\lineskiplimit| as well as |\baselineskip|. If
you say `|\openup2\jot|', the lines are spread apart 2 extra units, where
plain \TeX\ opens things up in units of $3\pt$. Since |$$...$$| acts as a
^{group}, the effect of\/ |\openup| will disappear when the display is
finished. Any \<dimen> can follow |\openup|, but it's customary to express the
amount symbolically in terms of a ^|\jot| instead of using absolute units;
^^{generic coding}
then your manuscript can be used with a variety of different formats.

\ddanger Plain \TeX's ^|\displaylines|, ^|\eqalignno|, and ^|\leqalignno|
macros begin with `|\openup1\jot|'. If you don't want the lines to be
opened up, you can cancel this by saying, e.g.,
`|$$\openup-1\jot \eqalignno{...}$$|', because |\openup| has a cumulative
effect.

\ddanger Suppose that you have decided to make a homegrown display having
the general form `|$$\openup1\jot \halign{...}$$|'; and for convenience,
let's suppose that the normal conventions of plain \TeX\ are in force, so
that |\jot=3pt| and |\baselineskip=12pt|. Then the |\openup| macro
changes the baselineskip distance to $15\pt$. It follows that the baseline
of the text line that immediately precedes the display will be $15\pt$
above the topmost baseline of the display, plus the ^|\abovedisplayskip|.
But when the paragraph resumes, its next baseline will be only $12\pt$
below the bottom baseline of the display, plus the ^|\belowdisplayskip|,
because the |\baselineskip| parameter will have reverted to its normal
value. The |\eqalignno| and |\displaylines| macros say
`|\noalign{\vskip|$-d$|}|' before their first lines, where $d$ is
the net amount of opening-up, in order to compensate for this difference.

\subsection Long formulas. Our discussion of mathematics typing is almost
complete; we need to deal with just one more problem: What should be
done when a formula is so long that it doesn't fit on a single line?

For example, suppose that you encounter the equation
$$\hfuzz=20pt % overfull box tolerated here
\sigma(2^{34}-1,2^{35},1)=
  -3+(2^{34}-1)/2^{35}+2^{35}\!/(2^{34}-1)+7/2^{35}(2^{34}-1)
  -\sigma(2^{35},2^{34}-1,1).$$ % from v2, 1st ed, p76
You'll have to break it up somehow; \TeX\ has done its best to squeeze
everything together by shrinking the spaces next to the $+$ and~$-$ signs
to zero, but still the line has come out overfull.

Let's try to break that equation just before the `$+7$'. One common way to
do this is to type
\begintt
$$\eqalign{\sigma(2^{34}-1,2^{35},1)
&=-3+(2^{34}-1)/2^{35}+2^{35}\!/(2^{34}-1)\cr
&\qquad+7/2^{35}(2^{34}-1)-\sigma(2^{35},2^{34}-1,1).\cr}$$
\endtt
which yields
$$\eqalign{\sigma(2^{34}-1,2^{35},1)
  &=-3+(2^{34}-1)/2^{35}+2^{35}\!/(2^{34}-1)\cr
  &\qquad+7/2^{35}(2^{34}-1)-\sigma(2^{35},2^{34}-1,1).\cr}$$
The idea is to treat a long one-line formula as a two-line formula,
using |\qquad| on the second line so that the second part of the formula
appears well to the right of the $=$~sign on the first line.

\exercise Explain how to deal with the following display. % v2 p107
$$\eqalignno{x_nu_1+\cdots+x_{n+t-1}u_t
   &=x_nu_1+(ax_n+c)u_2+\cdots\cr
   &\qquad+\bigl(a^{t-1}x_n+c(a^{t-2}+\cdots+1)\bigr)u_t\cr
   &=(u_1+au_2+\cdots+a^{t-1}u_t)x_n+h(u_1,\ldots,u_t).
     \quad&(47)\cr}$$
\answer |$$\eqalignno{x_nu_1+\cdots+x_{n+t-1}u_t|\parbreak
        |   &=x_nu_1+(ax_n+c)u_2+\cdots\cr|\parbreak
        |   &\qquad+\bigl(a^{t-1}x_n+c(a^{t-2}+\cdots+1)\bigr)u_t\cr|\parbreak
        |   &=(u_1+au_2+\cdots+a^{t-1}u_t)x_n+h(u_1,\ldots,u_t).|\parbreak
        |     \quad&(47)\cr}$$|\par\noindent
You weren't expected to insert the `|\quad|' on the last line; such
refinements usually can't be anticipated until you see the first proofs.
But without that |\quad| the `(47)' would occur half a quad closer to the
formula.

\danger It's quite an art to decide how to ^{break long displayed formulas}
into several lines; \TeX\ never attempts to break them, because no set of
rules is really adequate. The author of a mathematical manuscript is
generally the best judge of what to do, since break positions depend on
subtle factors of mathematical exposition. For example, it is often
desirable to emphasize some of the symmetry or other structure that
underlies a formula, and such things require a solid understanding of
exactly what is going on in that formula.

\begingroup\ninepoint
\danger Nevertheless, it is possible to state a few rules of thumb about
how to deal with long formulas in displays, since there are some
principles that the best mathematical typesetters tend to follow:\enddanger

\smallskip
\textindent{a)}Although formulas within a paragraph always break {\sl after\/}
binary operations and relations, displayed formulas always break {\sl before\/}
binary operations and relations. Thus, we didn't end the first line of
our $\sigma(\,\ldots\,)$ example with `|(2^{34}-1)+|'; we ended it with
`|(2^{34}-1)|' and began the second line with `|+|'.

\smallskip
\textindent{b)}When an equation is broken before a binary operation, the second
line should start at least two quads to the right of where the innermost
subformula containing that binary operation begins on the first line.
For example, if you wish to break
\begindisplay
|$$\sum_{0<k<n}\left(|\<formula$_1$>|+|\<formula$_2$>|\right)$$|
\enddisplay
at the plus sign between \<formula$_1$> and \<formula$_2$>, it is almost
mandatory to have the plus sign on the second line appear somewhat to the
right of the large left parenthesis that corresponds to `|\left(|'.

\endgroup

\danger In the example just considered, special care is needed to break the
formula into two lines, because ^|\left| and ^|\right| delimiters cannot be
used in isolation; you can't have only |\left| in one line of a formula
and only |\right| in the second. Furthermore, you'll want the two delimiters
to be of the same size, even though they occur in different lines. The best
solution is usually to choose the delimiter size yourself; for example,
you could type
\begindisplay
|$$\eqalign{\sum_{0<k<n}\biggl(&|\<formula$_1$>|\cr|\cr
|             &\qquad{}+|\<formula$_2$>|\biggr)\cr}$$|\cr
\enddisplay
if\/ ^|\bigg| delimiters are best. Notice that the |&|~markers don't occur
at $=$~signs in this example, they just mark a point of alignment.

\danger There's another way to break long formulas, sometimes called the
{\sl^{two-line}\/} form. The idea is to put the first part of the formula
almost ^{flush left}, and to put the second part almost ^{flush right},
where ``almost flush'' means ``one quad away.'' Thus, the two-line form of
the long $\sigma(\,\ldots\,)$ equation considered earlier is
$$\displaylines{\quad\sigma(2^{34}-1,2^{35},1)
  =-3+(2^{34}-1)/2^{35}+2^{35}\!/(2^{34}-1)\hfill\cr
\hfill{}+7/2^{35}(2^{34}-1)-\sigma(2^{35},2^{34}-1,1).\quad\cr}$$
It isn't difficult to get this two-line effect with ^|\displaylines|:
\begintt
$$\displaylines{\quad\sigma(2^{34}-1,2^{35},1)
  =-3+(2^{34}-1)/2^{35}+2^{35}\!/(2^{34}-1)\hfill\cr
\hfill{}+7/2^{35}(2^{34}-1)-\sigma(2^{35},2^{34}-1,1).\quad\cr}$$
\endtt
An extra `|{}|' was typed on the second line here so that \TeX\ would know
that the `|+|' is a binary operation.  The two-line form is especially
recommended for equations that have a long left-hand side; in that case the
break generally comes just before the~$=$~sign.

\dangerexercise Typeset the following display:
$$\displaylines{\quad\sum_{1\le j\le n}{1\over
    (x_j-x_1)\ldots(x_j-x_{j-1})(x-x_j)(x_j-x_{j+1})
    \ldots(x_j-x_n)}\hfill\cr
  \hfill={1\over(x-x_1)\ldots(x-x_n)}.\quad(27)\cr}$$ % v2 p80
\answer |$$\displaylines{\quad\sum_{1\le j\le n}{1\over|\parbreak
        |    (x_j-x_1)\ldots(x_j-x_{j-1})(x-x_j)(x_j-x_{j+1})|\parbreak
        |    \ldots(x_j-x_n)}\hfill\cr|\parbreak
        |  \hfill={1\over(x-x_1)\ldots(x-x_n)}.\quad(27)\cr}$$|

\ddangerexercise If it is necessary to typeset a huge fraction like
^^{fraction, huge}
$$\def\\#1;{(#1;q^2)_\infty}
q^{{1\over2}n(n+1)}\\ea;\\eq/a;\\caq/e;\\cq^2\!/ae;
\over(e;q)_\infty(cq/e;q)_\infty$$
in a single narrow column, you might have to break up the numerator and
resort to
\begindisplay
$\displaystyle{\def\\#1;{(#1;q^2)_\infty}
\displaystyle{q^{{1\over2}n(n+1)}\\ea;\\eq/a;\qquad\atop
  \hfill\\caq/e;\\cq^2\!/ae;}
\over(e;q)_\infty(cq/e;q)_\infty}$
\enddisplay
How would you specify the latter fraction to \TeX?
% cf SIAM J Math Anal 7 (1976) p333; even longer ones appear on p334
\answer |$$\def\\#1;{(#1;q^2)_\infty} % to save typing|\parbreak
        |\displaystyle{q^{{1\over2}n(n+1)}\\ea;\\eq/a;\qquad\atop|\parbreak
        |  \hfill\\caq/e;\\cq^2\!/ae;}|\parbreak
        |\over(e;q)_\infty(cq/e;q)_\infty$$|

\endchapter

When a formula is too long for the page-width
and has to be broken into successive lines
(and we are now, of course, speaking of displayed formulae),
it should be broken, if possible, at the end of a natural `phrase';
if, for example, it is a much-bracketed formula,
it should be broken at the end of one of the major brackets
and not at an inner symbol.
This natural phrasing (as in music or speech)
makes for intelligibility between writer and reader
and should not be left to the compositor.
An author, when he finds himself writing a longish formula,
should indicate a convenient point of fracture in case of need.
\author ^{CHAUNDY}, ^{BARRETT}, and ^{BATEY}, %
 {\sl The Printing of Mathematics\/} (1954) % p38

\bigskip

Some authors use display with discretion,
some run even extremely long, complicated equations into the text,
while others tend to display every equation in the paper.
The tendency to overdisplay is probably more predominant
than the tendency to underdisplay;
for this reason it is possible for the copy editor to shorten
(and even improve) papers by running displayed material into text. $\ldots$
On the other hand, there are occasions when the copy editor needs
to suggest the display of complicated expressions that have been run into text,
particularly when it would involve a bad break at the end of a text line.
\author ELLEN ^{SWANSON}, {\sl Mathematics into Type\/} (1971) % p41

\eject
